\section{FC-GNN: Fuzzy-Conformal Graph Neural Networks}
\label{sec:method}

Figure~\ref{fig:architecture} shows the FC-GNN architecture.
The model comprises three stages: (i) $L$ \emph{fuzzy message-passing layers}
that produce a per-node fuzzy membership distribution $\bm{\mu}_v \in [0,1]^C$,
(ii) an MLP output head, and (iii) post-hoc \emph{Fuzzy Mondrian Conformal
Prediction} calibration that wraps the trained model without retraining.

\begin{figure*}[t]
  \centering
  \includegraphics[width=\textwidth]{architecture}
  \caption{FC-GNN architecture.  Left: $L$ Fuzzy Message-Passing Layers
           (FMPL) compute Gaussian membership functions, product T-norm
           rule-firing strengths, and centroid defuzzification.
           Centre: MLP head produces fuzzy class memberships $\bm{\mu}_v$.
           Right: Fuzzy Mondrian CP partitions calibration nodes via Louvain
           clustering, computes Sugeno nonconformity scores, and returns a
           calibrated prediction set $\Calpha(v)$.}
  \label{fig:architecture}
\end{figure*}

\subsection{Fuzzy Message-Passing Layer (FMPL)}
\label{sec:fmpl}

Let $\bm{h}_v^{(\ell)} \in \R^d$ be the embedding of node $v$ at layer
$\ell$, with $\bm{h}_v^{(0)} = W_0 \bm{x}_v$.
The FMPL replaces the standard mean/sum aggregation with a fuzzy-rule-based
combiner operating as follows.

\paragraph{Gaussian membership functions.}
For each rule $k \in \{1, \ldots, K\}$, define a learnable Gaussian
membership function over the $d$-dimensional embedding space:
\[
  \mu_k(\bm{h}) = \exp\!\Bigl(
    -\tfrac{1}{2}\bigl\|\text{diag}(\bm{\sigma}_k)^{-1}(\bm{h}-\bm{c}_k)
    \bigr\|^2\Bigr),
\]
where centres $\bm{c}_k \in \R^d$ and log-widths
$\log\bm{\sigma}_k \in \R^d$ are learnable parameters.

\paragraph{Product T-norm rule firing.}
For node $v$ at layer $\ell$, the firing strength of rule $k$ is
\[
  \tau_k(v) = \prod_{u \in \calN(v)} \mu_k\bigl(\bm{h}_u^{(\ell)}\bigr),
\]
using the product T-norm, which is differentiable and suitable for
gradient-based learning.

\paragraph{Centroid defuzzification.}
Each rule $k$ has a consequent weight vector $\bm{w}_k \in \R^{d'}$.
The updated embedding is computed via Mamdani-style defuzzification:
\[
  \bm{h}_v^{(\ell+1)}
  = \frac{\sum_{k=1}^K \tau_k(v)\, \bm{w}_k}{\sum_{k=1}^K \tau_k(v)}
    + W_\text{res}\,\bm{h}_v^{(\ell)},
\]
where the residual term $W_\text{res}\bm{h}_v^{(\ell)}$ stabilises
gradient flow across layers.
An $\ell_1$ regularisation penalty on $\tau_k$ prunes inactive rules during
training.

\paragraph{MLP output head.}
After $L$ FMPL layers, an MLP with ReLU activation and dropout maps
$\bm{h}_v^{(L)}$ to logits, and a final softmax layer produces the
\emph{fuzzy class membership distribution}
$\bm{\mu}_v = (\mu_v^1, \ldots, \mu_v^C) \in [0,1]^C$,
$\sum_c \mu_v^c = 1$.

\subsection{Fuzzy Mondrian Conformal Prediction (FMCP)}
\label{sec:fmcp}

FMCP is applied post-hoc to a trained FC-GNN and proceeds in two phases:
calibration and inference.

\subsubsection{Sugeno Nonconformity Score}

For calibration node $v$ with true label $y_v$, define the
\emph{fuzzy nonconformity score}
\[
  s_v = 1 - \mathcal{S}_\lambda(\bm{\mu}_v, y_v),
\]
where $\mathcal{S}_\lambda(\bm{\mu}_v, y_v)$ is the Sugeno integral of the
membership distribution $\bm{\mu}_v$ with respect to the $\lambda$-fuzzy
measure, evaluated at the hypothesis ``the true class is $y_v$''.
Concretely, letting $\mu_{(1)} \ge \mu_{(2)} \ge \cdots \ge \mu_{(C)}$
be the sorted memberships and $g_{\lambda,k}$ the cumulative $\lambda$-measure
over the top-$k$ classes:
\[
  \mathcal{S}_\lambda(\bm{\mu}_v, y_v)
  = \Bigl(\max_{k=1}^C \min\bigl(\mu_{(k)},\, g_{\lambda,k}\bigr)\Bigr)
    \cdot \frac{\mu_v^{y_v}}{\max_c \mu_v^c + \varepsilon}.
\]
The scalar $\lambda$ is a learned parameter of the Sugeno module.
The second factor down-weights the integral by the relative membership of
the true class, so that $s_v$ is small when $\mu_v^{y_v}$ is large
(high confidence, correctly predicted) and large otherwise.

\subsubsection{Community Partition and Per-Community Quantiles}

We run Louvain community detection \citep{blondel08} on the full graph to
obtain a partition $\{Q_1, \ldots, Q_M\}$ of the node set.
Calibration nodes $V_\text{cal}$ are assigned to their community, and for
each community $Q_m$ the quantile
\[
  \hat{q}_{Q_m}
  = \hat{Q}_{1-\alpha}\bigl(\{s_v : v \in V_\text{cal} \cap Q_m\}\bigr)
\]
is computed at the $\lceil(|V_\text{cal} \cap Q_m|+1)(1-\alpha)\rceil$-th
order statistic.
Communities smaller than a minimum size threshold fall back to the global
quantile to maintain coverage.

\subsubsection{Prediction Set Construction}

At test time, node $v$ is assigned to its Louvain community $Q_m$ (or the
nearest community by feature-centroid distance for inductive nodes).
The fuzzy prediction set is
\[
  \Calpha(v)
  = \bigl\{c \in \{1,\ldots,C\} :
    1 - \mathcal{S}_\lambda(\bm{\mu}_v, c) \le \hat{q}_{Q_m}\bigr\}.
\]

\subsection{Coverage Theorem}
\label{sec:theorem}

We now state the formal coverage guarantee for FC-GNN.

\begin{definition}[Fuzzy Permutation Invariance]
\label{def:fpi}
Let $\pi$ be any permutation of $V$ preserving the graph structure
(i.e., $(u,v) \in \calE \Leftrightarrow (\pi(u),\pi(v)) \in \calE$).
The FC-GNN is \emph{fuzzy permutation-invariant} if the joint distribution
of Sugeno scores $(s_v)_{v \in V}$ is invariant under $\pi$ for all such $\pi$.
\end{definition}

\begin{assumption}[Consistent Community Assignment]
\label{ass:community}
For any node $v$ in the test set, the Louvain community assignment is
consistent with the assignment of its calibration neighbours with probability
at least $1 - \delta_0$ for a known $\delta_0 > 0$.
\end{assumption}

\begin{theorem}[FC-GNN Coverage]
\label{thm:coverage}
Let $G$ satisfy Definition~\ref{def:fpi} and Assumption~\ref{ass:community}.
For any test node $v$ assigned to community $Q_m$ at inference time:
\[
  \Prob\bigl(y_v \in \Calpha(v)\bigr) \ge 1 - \alpha.
\]
Furthermore, under Assumption~\ref{ass:community}:
\[
  \Prob\bigl(y_v \in \Calpha(v) \mid v \in Q_m\bigr)
  \ge 1 - \alpha - \delta(|Q_m|),
\]
where $\delta(|Q_m|) = O(|Q_m|^{-1/2})$ is a finite-sample slack term
arising from a Hoeffding inequality applied to the empirical quantile.
\end{theorem}

\begin{proof}[Proof sketch]
\textbf{Marginal coverage.}
By fuzzy permutation invariance (Definition~\ref{def:fpi}), the Sugeno
scores $(s_v)_{v \in V_\text{cal} \cup \{v_\text{test}\}}$ are jointly
exchangeable.
The standard CP exchangeability argument \citep{angelopoulos21} then gives
$\Prob(s_{v_\text{test}} \le \hat{q}) \ge 1-\alpha$, i.e.,
$\Prob(y_{v_\text{test}} \in \Calpha(v_\text{test})) \ge 1-\alpha$.

\textbf{Conditional coverage.}
Within community $Q_m$, the calibration scores $\{s_v : v \in Q_m\}$
are exchangeable by the same invariance argument.
The empirical quantile $\hat{q}_{Q_m}$ satisfies
$\Prob(s_v \le \hat{q}_{Q_m}) \ge 1-\alpha$ for any
test node $v$ exchangeable with the community calibration nodes.
By Hoeffding's inequality applied to the empirical c.d.f.\ of
$n_m = |Q_m \cap V_\text{cal}|$ i.i.d.\ Bernoulli trials,
the deviation of the empirical quantile from the true quantile is bounded by
$\delta(n_m) = \sqrt{\log(2/\delta_0) / (2n_m)} = O(n_m^{-1/2})$,
completing the proof.
\end{proof}

\begin{remark}
The slack $\delta(|Q_m|)$ vanishes as community size grows, matching the
behaviour of Mondrian CP in standard settings \citep{rrgnn}.
In practice, communities with fewer than five nodes fall back to the
global quantile, sacrificing conditional for marginal coverage in
data-sparse regimes.
\end{remark}

\subsection{Training Procedure}

Algorithm~\ref{alg:training} summarises the full FC-GNN pipeline.

\begin{algorithm}[t]
\caption{FC-GNN Training and Calibration}
\label{alg:training}
\begin{algorithmic}[1]
\Require Graph $G$, labels $Y_\text{train}$, calibration $V_\text{cal}$,
         level $\alpha$
\Ensure Trained FC-GNN, community quantiles $\{\hat{q}_{Q_m}\}$
\State \textbf{Phase 1 — Fuzzy GNN Training}
\State Initialise membership centres $\bm{c}_k$ and log-widths
       $\log\bm{\sigma}_k$ randomly
\For{each epoch}
  \State Forward pass: compute $\bm{\mu}_v$ for all $v \in V_\text{train}$
  \State $\mathcal{L} = \mathcal{L}_\text{CE} +
         \lambda_1 \|\bm{\tau}\|_1$
         \Comment{cross-entropy + rule sparsity}
  \State Backward pass; update all parameters via Adam
\EndFor
\State \textbf{Phase 2 — FMCP Calibration}
\State Forward pass on $V_\text{cal}$: obtain $\bm{\mu}_v$, $\forall v$
\State Run Louvain on $G$ to get communities $\{Q_1, \ldots, Q_M\}$
\For{each community $Q_m$}
  \State Compute $s_v = 1 - \mathcal{S}_\lambda(\bm{\mu}_v, y_v)$,
         $\forall v \in Q_m \cap V_\text{cal}$
  \State $\hat{q}_{Q_m} \leftarrow \lceil(|Q_m \cap V_\text{cal}|+1)(1-\alpha)\rceil$-th
         order statistic of $\{s_v\}$
\EndFor
\State \textbf{Phase 3 — Inference}
\For{each test node $v$}
  \State Assign $v$ to community $Q_m$ by Louvain / centroid proximity
  \State $\Calpha(v) \leftarrow \{c : 1 - \mathcal{S}_\lambda(\bm{\mu}_v, c)
         \le \hat{q}_{Q_m}\}$
  \State Extract top-$K$ fired rules for IF-THEN explanation
\EndFor
\end{algorithmic}
\end{algorithm}
